
\begin{filecontents*}{references.bib}
@misc{Frankl1979,
  author = {P{\'e}ter Frankl},
  title = {Union-closed sets conjecture},
  year = {1979},
  note = {Problem statement; see also D. B. West's Open Problems in Graph Theory and Combinatorics},
  url = {https://dwest.web.illinois.edu/openp/unionclos.html}
}
@article{BruhnSchaudt2015,
  author = {Henning Bruhn and Oliver Schaudt},
  title = {The journey of the union-closed sets conjecture},
  journal = {Graphs and Combinatorics},
  volume = {31},
  number = {6},
  pages = {2043--2074},
  year = {2015},
  doi = {10.1007/s00373-014-1515-0}
}
@misc{Gilmer2022,
  author = {Justin Gilmer},
  title = {A constant lower bound for the union-closed sets conjecture},
  year = {2022},
  eprint = {2211.09055},
  archivePrefix = {arXiv},
  primaryClass = {math.CO}
}
@article{AlweissHuangSellke2024,
  author = {Ryan Alweiss and Brice Huang and Mark Sellke},
  title = {Improved Lower Bound for Frankl's Union-Closed Sets Conjecture},
  journal = {The Electronic Journal of Combinatorics},
  volume = {31},
  number = {3},
  pages = {P3.35},
  year = {2024},
  doi = {10.37236/12232}
}
@misc{ChaseLovett2022,
  author = {Zachary Chase and Shachar Lovett},
  title = {Approximate union closed conjecture},
  year = {2022},
  eprint = {2211.11689},
  archivePrefix = {arXiv},
  primaryClass = {math.CO}
}
@misc{Sawin2022,
  author = {Will Sawin},
  title = {An improved lower bound for the union-closed sets conjecture},
  year = {2022},
  eprint = {2211.11504},
  archivePrefix = {arXiv},
  primaryClass = {math.CO}
}
@misc{Liu2023,
  author = {Jingbo Liu},
  title = {Improving the Lower Bound for the Union-closed Sets Conjecture via Conditionally IID Coupling},
  year = {2023},
  eprint = {2306.08824},
  archivePrefix = {arXiv},
  primaryClass = {math.CO}
}
@misc{Cambie2022,
  author = {Stijn Cambie},
  title = {Better bounds for the union-closed sets conjecture using the entropy approach},
  year = {2022},
  eprint = {2212.12500},
  archivePrefix = {arXiv},
  primaryClass = {math.CO}
}
@article{Yu2023,
  author = {Lei Yu},
  title = {Dimension-Free Bounds for the Union-Closed Sets Conjecture},
  journal = {Entropy},
  volume = {25},
  number = {5},
  pages = {767},
  year = {2023},
  doi = {10.3390/e25050767}
}
@misc{Pebody2022,
  author = {L. Pebody},
  title = {Extension of a Method of Gilmer},
  year = {2022},
  eprint = {2211.13139},
  archivePrefix = {arXiv},
  primaryClass = {math.CO}
}
@inproceedings{Boppana1985,
  author = {R. B. Boppana},
  title = {Amplification of Probabilistic Boolean Formulas},
  booktitle = {26th Annual Symposium on Foundations of Computer Science (SFCS 1985)},
  pages = {20--29},
  publisher = {IEEE},
  year = {1985}
}
@article{PulajWood2024,
  author = {Jonad Pulaj and Kenan Wood},
  title = {Local Configurations in Union-Closed Families},
  journal = {Experimental Mathematics},
  year = {2024},
  note = {To appear; arXiv:2301.01331}
}
\end{filecontents*}

\documentclass[12pt]{article}
\usepackage[margin=1in]{geometry}
\usepackage{amsmath, amsthm, amssymb}
\usepackage{hyperref}
\usepackage{booktabs}
\usepackage{enumitem}
\usepackage{longtable}
\usepackage{mathtools}
\allowdisplaybreaks

\newtheorem{theorem}{Theorem}
\newtheorem{lemma}[theorem]{Lemma}
\newtheorem{proposition}[theorem]{Proposition}
\newtheorem{corollary}[theorem]{Corollary}
\newtheorem{claim}[theorem]{Claim}
\theoremstyle{definition}
\newtheorem{definition}[theorem]{Definition}
\theoremstyle{remark}
\newtheorem{remark}[theorem]{Remark}

\numberwithin{equation}{section}

\newcommand{\eps}{\varepsilon}
\newcommand{\R}{\mathbb{R}}
\newcommand{\Fcal}{\mathcal{F}}
\newcommand{\E}{\mathbb{E}}
\newcommand{\Prob}{\mathbb{P}}
\newcommand{\one}{\mathbf{1}}
\newcommand{\bstar}{\beta_*}
\newcommand{\mstar}{m_*}
\newcommand{\xstar}{x_*}
\newcommand{\alphastar}{\alpha_*}
\newcommand{\cstar}{c'}

\title{A Certified Bivariate Entropy Inequality Related to Liu's Entropy Bound}
\author{K\aa re}
\date{}

\begin{document}
\maketitle

\begin{abstract}
We prove, with certified computation, a bivariate entropy inequality that is
suggested by Liu's conditionally independent coupling method for the
union-closed sets conjecture.  The proof combines a Boppana-style Rolle method
on the diagonal, a degree-27 Bernstein positivity certificate, and a finite
Taylor-model interval-arithmetic verification of the off-diagonal
critical-point structure.  This version records the verified bivariate
inequality and its computational certificates.  It does not claim the
union-closed corollary: the previously proposed reduction from the pointwise
bivariate inequality to Liu's full conditionally i.i.d. variational problem
implicitly used an independence assumption that is not available in the
conditionally i.i.d. coupling.  The remaining task is to connect the certified
bivariate inequality to the correct distributional optimization problem, or to
replace it by a stronger distributional certificate.  The code implementing all
certificates is available at the repository listed in the appendix.
\end{abstract}

\paragraph{Status note.}
An earlier draft claimed that the bivariate inequality proved here implied the
constant \(0.382709087918735\ldots\) for Frankl's union-closed sets conjecture.
Jingbo Liu pointed out that the entropy argument in that draft implicitly
identified two different distributions of conditional probabilities.  In Liu's
conditionally i.i.d. framework, the independent entropy term and the
conditionally independent entropy term are integrated against different
measures.  The present version therefore withdraws the union-closed corollary
and keeps only the certified bivariate inequality and the precise statement of
the remaining distributional gap.

\section{Introduction}

A finite family \(\Fcal\) of sets is \emph{union-closed} if
\begin{equation}
  A,B\in\Fcal \quad\Longrightarrow\quad A\cup B\in\Fcal.
\end{equation}
Frankl's union-closed sets conjecture, formulated in 1979, asserts that if
\(\Fcal\) is finite, union-closed, and not equal to \(\{\varnothing\}\), then
some element belongs to at least half of the sets in \(\Fcal\)
\cite{Frankl1979,BruhnSchaudt2015}.  Despite its elementary formulation, the
conjecture has remained open for more than four decades.  It is one of the
central extremal problems concerning finite set systems.

A major breakthrough was made by Gilmer \cite{Gilmer2022}, who introduced an
entropy method and proved the first absolute constant lower bound.  In
Gilmer's argument, one samples independent random sets \(A,B\) from
\(\Fcal\), compares \(H(A\cup B)\) with \(H(A)\), and uses union-closedness to
force a popular element.  Gilmer's initial constant was \(0.01\).  This was the
first proof that there is an absolute \(c>0\) such that every finite nonempty
union-closed family different from \(\{\varnothing\}\) has an element appearing
in at least a \(c\)-fraction of its members.

Shortly afterward, the constant was improved to
\begin{equation}
  \frac{3-\sqrt 5}{2}=0.3819660112\ldots
\end{equation}
by several works, including Alweiss--Huang--Sellke
\cite{AlweissHuangSellke2024}, Chase--Lovett \cite{ChaseLovett2022}, Sawin
\cite{Sawin2022}, Cambie \cite{Cambie2022}, Yu \cite{Yu2023}, and Pebody
\cite{Pebody2022}.  A central analytic ingredient in this line of work is the
one-variable entropy inequality
\begin{equation}\label{eq:intro_boppana}
  h(x^2)\ge \varphi xh(x),
  \qquad
  0\le x\le1,
  \qquad
  \varphi=\frac{1+\sqrt5}{2},
\end{equation}
which goes back to Boppana's work on amplification of probabilistic Boolean
formulas \cite{Boppana1985}.  Chase and Lovett also showed that the constant
\((3-\sqrt5)/2\) is optimal for an approximate union-closed formulation,
thereby identifying a barrier for the plain independent two-sample entropy
method.

Liu \cite{Liu2023} introduced a further refinement based on conditionally
independent couplings.  Instead of taking two independent samples, one allows
the two samples to be independent only after conditioning on an auxiliary
random variable.  This produces a richer class of coordinatewise couplings and
leads to the numerical constant
\begin{equation}
  0.3827090879\ldots,
\end{equation}
slightly above the golden-ratio threshold.  Liu proved analytically a strict
improvement over the previous \(0.38234\)-level constants, but the explicit
value \(0.38271\ldots\) in his paper depends on additional numerically verified
hypotheses.  More precisely, Liu's explicit value requires two extra
certificates: first, a positive-semidefinite condition for an infinite kernel
matrix associated with the conditionally independent protocol; second, a
global optimizer-structure hypothesis for a nine-dimensional variational
problem.  The optimizer is observed numerically to have a two-point boundary
structure, but Liu does not prove this structure analytically.

The present paper studies the bivariate inequality that one obtains by taking
the apparent two-point boundary optimizer suggested by Liu's computation and
writing down the corresponding pointwise entropy certificate.  Let
\begin{equation}
  I(x,y)=xy(1+(1-x)(1-y)).
\end{equation}
We prove that, for explicit constants \(\mstar\) and \(\bstar\),
\begin{equation}\label{eq:intro_main_ineq}
  \mstar\bigl((1-\bstar)h(xy)+\bstar h(I(x,y))\bigr)
  \ge
  \frac{x h(y)+y h(x)}2
  \qquad (0\le x,y\le1).
\end{equation}
Here
\begin{equation}
  \mstar=0.617290912081265\ldots,
  \qquad
  1-\mstar=0.382709087918735\ldots .
\end{equation}
The inequality \eqref{eq:intro_main_ineq} is a certified bivariate entropy
inequality related to Liu's proposed constant.  It is not, by itself, a proof
of the corresponding union-closed lower bound.  The reason is distributional:
in Liu's variational problem the independent term and the conditionally
independent term are averaged with respect to different measures, whereas a
pointwise inequality such as \eqref{eq:intro_main_ineq} controls both terms on
the same pair \((x,y)\).  A further distributional argument would be needed to
recover the full union-closed consequence.

The proof of \eqref{eq:intro_main_ineq} has three parts.  On the boundary of
the square \([0,1]^2\), the inequality is immediate from \(m_*>1/2\).  On the
diagonal \(x=y\), we write \(M(x)=G(x,x)\) and prove \(M(x)\ge0\) by a
fifth-derivative sign argument: a symbolic computation reduces the sign of
\(M^{(5)}\) to the positivity of a degree-\(27\) polynomial \(P\), and we
certify \(P>0\) on \([0,1]\) by verifying the positivity of all \(28\) of its
Bernstein coefficients.  Off the diagonal, any negative interior minimum would
be an off-diagonal critical point.  Passing to variables
\[
  u=\frac{x+y}{2},
  \qquad
  v=\frac{x-y}{2},
\]
we analyze the branch \(F=\Gamma_v/v=0\).  Along this branch, the derivative of
\(\Gamma_u\) is controlled by the Schur-complement identity
\[
  \frac{d}{du}\Gamma_u(u,v(u))
  =
  \frac{\det D^2\Gamma(u,v(u))}{\Gamma_{vv}(u,v(u))}.
\]
A Taylor-model interval computation certifies the required determinant signs
on the two branch arms, while a separate interval-Newton and topological
argument certifies that there are no additional components of \(F=0\).  This
rules out all off-diagonal critical points.

Computer-assisted inequalities have already played a role in the recent
entropy approach to the union-closed conjecture.  In particular,
Alweiss--Huang--Sellke used computer verification in their treatment of the
one-variable inequality underlying the \((3-\sqrt5)/2\) bound
\cite{AlweissHuangSellke2024}.  The present paper uses certified computation
in the same spirit: all numerical claims are reduced to finite interval
statements with explicit margins.  The analytic part of the proof identifies
the correct bivariate inequality and reduces the off-diagonal problem to
branch isolation and determinant sign checks; the appendix records the finite
certificates needed to make those reductions rigorous.

\section{Definitions and constants}\label{sec:defs}

Throughout the paper we use natural logarithms.  Define the binary entropy function by
\begin{equation}\label{eq:hdef}
  h(t)=-t\log t-(1-t)\log(1-t),
  \qquad 0\le t\le 1,
\end{equation}
with the convention \(0\log0=0\).  Let
\begin{equation}\label{eq:Ldef}
  L(t)=h'(t)=\log\frac{1-t}{t}.
\end{equation}
Define
\begin{equation}\label{eq:Idef}
  I(x,y)=xy(1+(1-x)(1-y))=xy(2-x-y+xy),
\end{equation}
and
\begin{equation}\label{eq:Ddef}
  D(x)=I(x,x)=x^2(2-2x+x^2).
\end{equation}

Let \(\xstar\in(0,1)\) be the unique root of
\begin{equation}\label{eq:xstarpoly}
  x^4-2x^3+3x^2-1=0.
\end{equation}
Equivalently,
\begin{equation}\label{eq:xstarequiv}
  \xstar^2(2+(1-\xstar)^2)=1.
\end{equation}
Set
\begin{equation}\label{eq:pstar}
  p_* = \frac{h(\xstar)}{h(\xstar^2)},
\end{equation}
\begin{equation}\label{eq:mstar}
  \mstar=p_*\xstar,
  \qquad
  c'=1-\mstar.
\end{equation}
Numerically,
\begin{equation}\label{eq:mstar_numeric}
  \xstar=0.690787593924988\ldots,
  \qquad
  \mstar=0.617290912081265\ldots,
  \qquad
  c'=0.382709087918735\ldots .
\end{equation}

The parameter \(\bstar\) is defined by the stationarity condition
\begin{equation}\label{eq:betastar_stationary}
  \frac{d}{dx}\left[
  \mstar\bigl((1-\bstar)h(x^2)+\bstar h(D(x))\bigr)-xh(x)
  \right]_{x=\xstar}=0.
\end{equation}
Equivalently,
\begin{equation}\label{eq:betastar_formula}
  \bstar=
  \frac{
  \dfrac{h(\xstar^2)}{h(\xstar)}
  \left(h'(\xstar)+\dfrac{h(\xstar)}{\xstar}\right)
  -2\xstar h'(\xstar^2)
  }{
  D'(\xstar)h'(D(\xstar))-2\xstar h'(\xstar^2)
  }.
\end{equation}
Numerically,
\begin{equation}\label{eq:betastar_numeric}
  \bstar=0.100052559862893\ldots .
\end{equation}
We write
\begin{equation}\label{eq:alpha}
  \alphastar=1-\bstar.
\end{equation}

Define
\begin{equation}\label{eq:Gdef}
  G(x,y)=\mstar\bigl(\alphastar h(xy)+\bstar h(I(x,y))\bigr)-\frac{x h(y)+y h(x)}2.
\end{equation}
The main analytic result of the paper is the inequality
\begin{equation}\label{eq:main_bivar}
  G(x,y)\ge 0
  \qquad(0\le x,y\le 1).
\end{equation}

We record the equality identities at \((\xstar,\xstar)\).  From \eqref{eq:xstarpoly} and \eqref{eq:Ddef},
\begin{equation}\label{eq:Dxstar}
  D(\xstar)=1-\xstar^2.
\end{equation}
By entropy symmetry,
\begin{equation}\label{eq:h_sym_xstar}
  h(D(\xstar))=h(1-\xstar^2)=h(\xstar^2).
\end{equation}
Therefore
\begin{equation}\label{eq:Gxstarzero}
  G(\xstar,\xstar)=\mstar h(\xstar^2)-\xstar h(\xstar)=0,
\end{equation}
where the last equality is \eqref{eq:pstar} and \eqref{eq:mstar}.  Finally, \eqref{eq:betastar_stationary} says
\begin{equation}\label{eq:grad_xstar}
  \frac{d}{dx}G(x,x)\bigg|_{x=\xstar}=0.
\end{equation}
Since \(G\) is symmetric, \eqref{eq:grad_xstar} implies
\begin{equation}\label{eq:gradient_zero}
  \nabla G(\xstar,\xstar)=0.
\end{equation}

\section{Boundary cases}\label{sec:boundary}

We first prove nonnegativity on the boundary of the square.  If \(x=0\), then \(xy=0\), \(I(0,y)=0\), and \(xh(y)=0\), hence
\begin{equation}\label{eq:axis1}
  G(0,y)=0.
\end{equation}
Similarly,
\begin{equation}\label{eq:axis2}
  G(x,0)=0.
\end{equation}
If \(x=1\), then \(xy=y\) and \(I(1,y)=y\), so
\begin{equation}\label{eq:xone}
  G(1,y)=\left(\mstar-\frac12\right)h(y)\ge 0,
\end{equation}
because \(\mstar>1/2\).  By symmetry,
\begin{equation}\label{eq:yone}
  G(x,1)=\left(\mstar-\frac12\right)h(x)\ge 0.
\end{equation}
Thus \(G\ge0\) on \(\partial[0,1]^2\), with equality on the axes and at the corner \((1,1)\).

\section{The diagonal inequality}\label{sec:diagonal}

Define
\begin{equation}\label{eq:Mdef}
  M(x)=G(x,x).
\end{equation}
Using \(I(x,x)=D(x)\), we have
\begin{equation}\label{eq:Mformula}
  M(x)=\mstar\bigl(\alphastar h(x^2)+\bstar h(D(x))\bigr)-xh(x).
\end{equation}

\begin{lemma}[Diagonal nonnegativity]\label{lem:diagonal}
For every \(x\in[0,1]\),
\begin{equation}\label{eq:Mnonneg}
  M(x)\ge0.
\end{equation}
On \([\mstar,1]\), equality occurs only at \(x=\xstar\) and \(x=1\).  On \([0,1]\), equality occurs at \(x=0\), \(x=\xstar\), and \(x=1\).
\end{lemma}

\begin{proof}
A symbolic differentiation gives
\begin{equation}\label{eq:M5}
  M^{(5)}(x)=-\frac{2P(x)}{x^3(1-x)^4(1+x)^4(x^2-2x+2)^4(x^3-x^2+x+1)^4}.
\end{equation}
The denominator is strictly positive on \((0,1)\).  The polynomial \(P\) is
\begin{equation}\label{eq:Pdecomp}
  P(x)=P_0(x)+\mstar P_1(x)+\mstar\bstar P_2(x).
\end{equation}
The three polynomials are
\begin{align}\label{eq:P0}
P_0(x)=&\;3x^{27}-36x^{26}+200x^{25}-666x^{24}+1364x^{23}-1360x^{22}
-1096x^{21}+6288x^{20}\notag\\
&-9854x^{19}+4312x^{18}+11216x^{17}-23492x^{16}+15180x^{15}+11376x^{14}-29832x^{13}\notag\\
&+18608x^{12}+9251x^{11}-22292x^{10}+9976x^9+6670x^8-9176x^7+1408x^6\notag\\
&+2784x^5-1488x^4-400x^3+320x^2-32,
\end{align}
\begin{align}\label{eq:P1}
P_1(x)=&\;24x^{26}-288x^{25}+1860x^{24}-8208x^{23}+27296x^{22}-71472x^{21}+150320x^{20}\notag\\
&-254064x^{19}+338288x^{18}-333680x^{17}+197272x^{16}+27248x^{15}-209888x^{14}\notag\\
&+233136x^{13}-104048x^{12}-45904x^{11}+96120x^{10}-46640x^9-12924x^8\notag\\
&+25312x^7-6336x^6-5312x^5+3360x^4+640x^3-640x^2+64,
\end{align}
\begin{align}\label{eq:P2}
P_2(x)=&\;24x^{26}-168x^{25}+60x^{24}+2832x^{23}-14144x^{22}+41712x^{21}-99488x^{20}\notag\\
&+213456x^{19}-389072x^{18}+532880x^{17}-458888x^{16}+88048x^{15}+372160x^{14}\notag\\
&-587200x^{13}+435712x^{12}-83408x^{11}-223240x^{10}+324792x^9-210324x^8\notag\\
&+22720x^7+77632x^6-70096x^5+29760x^4-5760x^3-256x^2+192x+64.
\end{align}

The degree-27 Bernstein coefficients \(b_i\) of \(P\) on \([0,1]\) have certified lower bounds shown in Table \ref{tab:bernstein}.  Since the Bernstein basis is nonnegative on \([0,1]\), the table implies
\begin{equation}\label{eq:Ppositive}
  P(x)>0\qquad(0\le x\le1).
\end{equation}
Consequently,
\begin{equation}\label{eq:M5negative}
  M^{(5)}(x)<0\qquad(0<x<1).
\end{equation}

\begin{table}[ht]
\centering
\caption{Certified lower bounds for the Bernstein coefficients of \(P\).}
\label{tab:bernstein}
\begin{tabular}{cr@{\quad}cr}
\toprule
\(i\) & \multicolumn{1}{c}{lower bound} & \(i\) & \multicolumn{1}{c}{lower bound} \\
\midrule
0 & 11.45935667296576 & 14 & 4.610838565116360\\
1 & 11.898549817384106 & 15 & 4.964896105425193\\
2 & 12.078833739776622 & 16 & 5.668547978061419\\
3 & 11.876898948434980 & 17 & 6.798303977430921\\
4 & 11.307562491478990 & 18 & 8.465673305310538\\
5 & 10.464009676814568 & 19 & 10.840728674923076\\
6 & 9.465777139811940 & 20 & 14.191030026624640\\
7 & 8.423319403297192 & 21 & 18.945974671912115\\
8 & 7.421560841110654 & 22 & 25.804885986190340\\
9 & 6.518660405060565 & 23 & 35.924168951745940\\
10 & 5.753064074484635 & 24 & 51.256141180419560\\
11 & 5.152687277174934 & 25 & 75.205102681605100\\
12 & 4.742718854790539 & 26 & 114.019777429866760\\
13 & 4.550899799454681 & 27 & 180.158840956815370\\
\bottomrule
\end{tabular}
\end{table}

We first prove the claim on \([\mstar,1]\).  The following signs are certified by interval evaluation:
\begin{align}\label{eq:mstar_signs}
&M(\mstar)>0,\qquad M'(\mstar)<0,
\qquad M''(\mstar)>0,\notag\\
&M^{(3)}(\mstar)>0,
\qquad M^{(4)}(\mstar)<0.
\end{align}
Numerically these are
\begin{align}\label{eq:mstar_signs_num}
&M(\mstar)=0.000777106330304900\ldots,\notag\\
&M'(\mstar)=-0.0199407396705102\ldots,\notag\\
&M''(\mstar)=0.219173887997792\ldots,\notag\\
&M^{(3)}(\mstar)=1.55739338128487\ldots,\notag\\
&M^{(4)}(\mstar)=-4.41104974811261\ldots .
\end{align}
Moreover,
\begin{equation}\label{eq:stationary_M}
  M(\xstar)=0,
  \qquad
  M'(\xstar)=0,
\end{equation}
and interval evaluation gives
\begin{equation}\label{eq:M2xstar}
  M''(\xstar)=0.317062949668787\ldots>0.
\end{equation}

We also need the behavior at \(1\).  Put \(s=1-x\).  As \(s\downarrow0\),
\begin{align}\label{eq:endpoint_h}
  h(x^2)&=2s\log\frac1s+O(s),\notag\\
  h(D(x))&=2s\log\frac1s+O(s),\notag\\
  xh(x)&=s\log\frac1s+O(s).
\end{align}
Therefore
\begin{equation}\label{eq:M_endpoint_expand}
  M(x)=(2\mstar-1)s\log\frac1s+O(s).
\end{equation}
Since \(2\mstar-1>0\), differentiating with respect to \(x=1-s\) gives
\begin{align}\label{eq:M_deriv_endpoint}
  M'(1^-)&=-\infty,\notag\\
  M''(1^-)&=-\infty,\notag\\
  M^{(3)}(1^-)&=-\infty,\notag\\
  M^{(4)}(1^-)&=-\infty.
\end{align}

By \eqref{eq:M5negative}, \(M^{(4)}\) is strictly decreasing.  Since \(M^{(4)}(\mstar)<0\), we have
\begin{equation}\label{eq:M4neg}
  M^{(4)}(x)<0\qquad(\mstar\le x<1).
\end{equation}
Hence \(M^{(3)}\) is strictly decreasing.  From \(M^{(3)}(\mstar)>0\) and \(M^{(3)}(1^-)=-\infty\), there is a unique
\begin{equation}\label{eq:a_zero}
  a\in(\mstar,1)
\end{equation}
with \(M^{(3)}(a)=0\).  Thus \(M''\) increases on \((\mstar,a)\) and decreases on \((a,1)\).  Since \(M''(\mstar)>0\) and \(M''(1^-)=-\infty\), there is a unique
\begin{equation}\label{eq:b_zero}
  b\in(a,1)
\end{equation}
with \(M''(b)=0\).

Consequently \(M'\) increases on \((\mstar,b)\) and decreases on \((b,1)\).  We have \(M'(\mstar)<0\), \(M'(\xstar)=0\), and \(M''(\xstar)>0\), so \(\xstar\in(\mstar,b)\) and is the unique zero of \(M'\) in \((\mstar,b)\).  Also \(M'(b)>0\), while \(M'(1^-)=-\infty\); hence there is a unique
\begin{equation}\label{eq:d_zero}
  d\in(b,1)
\end{equation}
with \(M'(d)=0\).  The sign pattern is
\begin{equation}\label{eq:Mprime_pattern}
  M'<0\text{ on }(\mstar,\xstar),\qquad
  M'>0\text{ on }(\xstar,d),\qquad
  M'<0\text{ on }(d,1).
\end{equation}
Since \(M(\mstar)>0\), \(M(\xstar)=0\), and \(M(1)=0\), it follows that
\begin{equation}\label{eq:M_ge_mstar}
  M(x)\ge0\qquad(\mstar\le x\le1).
\end{equation}

It remains to treat \([0,\mstar]\).  As \(x\downarrow0\),
\begin{align}\label{eq:M_zero_expand}
  h(x^2)&=2x^2\log\frac1x+O(x^2),\notag\\
  h(D(x))&=4x^2\log\frac1x+O(x^2),\notag\\
  xh(x)&=x^2\log\frac1x+O(x^2).
\end{align}
Hence
\begin{equation}\label{eq:M_zero_leading}
  M(x)=\gamma x^2\log\frac1x+O(x^2),
  \qquad
  \gamma=2\mstar(1+\bstar)-1>0.
\end{equation}
Thus \(M(x)>0\) for all sufficiently small positive \(x\), \(M'(0^+)=0\), \(M''(0^+)=+\infty\), \(M^{(3)}(0^+)=-\infty\), and \(M^{(4)}(0^+)=+\infty\).

Since \(M^{(5)}<0\), \(M^{(4)}\) is strictly decreasing on \((0,\mstar)\).  From \(M^{(4)}(0^+)=+\infty\) and \(M^{(4)}(\mstar)<0\), there is a unique zero \(c\in(0,\mstar)\) of \(M^{(4)}\).  Thus \(M^{(3)}\) increases on \((0,c)\) and decreases on \((c,\mstar)\).  Since \(M^{(3)}(0^+)=-\infty\) and \(M^{(3)}(\mstar)>0\), \(M^{(3)}\) has exactly one zero \(a_0\in(0,c)\).  Therefore \(M''\) decreases on \((0,a_0)\) and increases on \((a_0,\mstar)\).

The interval evaluation
\begin{equation}\label{eq:M2_01}
  M''(0.1)<0
\end{equation}
combined with \(M''(0^+)=+\infty\) and \(M''(\mstar)>0\) implies that \(M''\) has exactly two zeros \(b_1<b_2\) in \((0,\mstar)\).  Thus \(M'\) increases, then decreases, then increases.  Since \(M'(0^+)=0\) and \(M'(x)>0\) for small \(x>0\), while \(M'(\mstar)<0\), the function \(M'\) has exactly one zero \(r\in(0,\mstar)\).  Hence \(M\) increases on \((0,r)\) and decreases on \((r,\mstar)\).  Since \(M(0)=0\) and \(M(\mstar)>0\), we get
\begin{equation}\label{eq:M_ge_zero_m}
  M(x)\ge0\qquad(0\le x\le\mstar).
\end{equation}
This completes the proof.
\end{proof}

\section{Off-diagonal critical points}\label{sec:offdiagonal}

Introduce
\begin{equation}\label{eq:uvdef}
  u=\frac{x+y}{2},
  \qquad
  v=\frac{x-y}{2}.
\end{equation}
Thus
\begin{equation}\label{eq:xy_uv}
  x=u+v,
  \qquad
  y=u-v.
\end{equation}
The feasible half-triangle is
\begin{equation}\label{eq:feasible_triangle}
  0<u<1,
  \qquad
  0<v<\min(u,1-u).
\end{equation}
Define
\begin{equation}\label{eq:Gamma_def}
  \Gamma(u,v)=G(u+v,u-v).
\end{equation}
Since \(G\) is symmetric, \(\Gamma\) is even in \(v\).  For \(v>0\), define
\begin{equation}\label{eq:Fdef}
  F(u,v)=\frac{\Gamma_v(u,v)}{v}.
\end{equation}
An off-diagonal critical point satisfies
\begin{equation}\label{eq:critical_conditions}
  F(u,v)=0,
  \qquad
  \Gamma_u(u,v)=0,
  \qquad
  v>0.
\end{equation}

Set
\begin{equation}\label{eq:abJ}
  a=u^2-v^2,
  \qquad
  b=2-2u+u^2-v^2,
  \qquad
  J=ab.
\end{equation}
Thus \(a=xy\), \(b=1+(1-x)(1-y)\), and \(J=I(x,y)\).  Let
\begin{equation}\label{eq:Delta_def}
  \Delta(u,v)=\det D^2\Gamma(u,v)=\Gamma_{uu}\Gamma_{vv}-\Gamma_{uv}^2.
\end{equation}
Along a regular branch \(v=v(u)\) of \(F=0\), one has \(\Gamma_v(u,v(u))=0\).  Differentiating gives
\begin{equation}\label{eq:vprime}
  v'(u)=-\frac{\Gamma_{uv}}{\Gamma_{vv}}.
\end{equation}
For
\begin{equation}\label{eq:Phi_def}
  \Phi(u)=\Gamma_u(u,v(u)),
\end{equation}
we therefore get
\begin{equation}\label{eq:schur}
  \Phi'(u)=\Gamma_{uu}-\frac{\Gamma_{uv}^2}{\Gamma_{vv}}=\frac{\Delta}{\Gamma_{vv}}.
\end{equation}
This identity is the monotonicity mechanism for the off-diagonal analysis.

\subsection{Endpoint positivity on the upper arc}

Put
\begin{equation}\label{eq:eps_s_coords}
  x=1-\eps,
  \qquad
  y=s.
\end{equation}
Let
\begin{equation}\label{eq:q_def}
  q=2\mstar-1=0.234581824162530\ldots .
\end{equation}
On the upper branch, \(F=0\) is equivalent to \(G_x=G_y\), hence
\begin{equation}\label{eq:Gamma_u_2Gy}
  \Gamma_u=G_x+G_y=2G_y.
\end{equation}
Since
\begin{equation}\label{eq:Gy_formula}
  G_y=m_*\bigl(\alphastar xL(xy)+\bstar I_y(x,y)L(I(x,y))\bigr)-\frac{h(x)+xL(y)}2,
\end{equation}
we have the exact identity
\begin{equation}\label{eq:Gamma_u_endpoint_exact}
  \Gamma_u=2m_*\bigl(\alphastar xL(xy)+\bstar I_y(x,y)L(I(x,y))\bigr)-h(x)-xL(y).
\end{equation}
Here
\begin{equation}\label{eq:I_y_endpoint}
  I_y(x,y)=x(2-x-2y+2xy).
\end{equation}

\begin{lemma}[Endpoint positivity]\label{lem:endpoint}
For every point of the upper arc with \(0<s=y\le s_1=0.04\),
\begin{equation}\label{eq:endpoint_positive}
  \Gamma_u\ge \frac{2\mstar-1}{2}\log\frac1s>0.
\end{equation}
\end{lemma}

\begin{proof}
The endpoint branch satisfies
\begin{equation}\label{eq:eps_s2}
  0<\eps(s)\le s^2\qquad(0<s\le0.04).
\end{equation}
Indeed, from the explicit formula for \(F=(G_x-G_y)/v\), one obtains
\begin{align}\label{eq:F_endpoint_bounds}
  &\lim_{\eps\downarrow0}F(\eps,s)=+\infty,\notag\\
  &F(s^2,s)\le -\frac{q}{2}\log\frac1s+0.003<0,
\end{align}
and
\begin{equation}\label{eq:Feps_negative}
  \partial_\eps F(\eps,s)\le -\frac{1}{4s}<0
  \qquad(0<\eps\le s^2,\,0<s\le0.04).
\end{equation}
Thus the unique endpoint solution satisfies \eqref{eq:eps_s2}.

Assume \(0<s\le s_1\) and \(0<\eps\le s^2\).  Then \(x=1-\eps\geq 1-s^2\).  Since \(xy\le s\) and \(L\) is decreasing,
\begin{equation}\label{eq:Lxy_bound}
  L(xy)\ge L(s).
\end{equation}
Moreover,
\begin{equation}\label{eq:Iy_lower}
  I_y=(1-\eps)(1+\eps-2\eps s)\geq 1-s^2,
\end{equation}
and
\begin{equation}\label{eq:I_upper_endpoint}
  I(x,y)\le s(1+s^2).
\end{equation}
Therefore
\begin{equation}\label{eq:LI_bound}
  L(I(x,y))\ge L(s(1+s^2)).
\end{equation}
Also \(h(x)=h(\eps)\le h(s^2)\) and \(-xL(y)\ge -L(s)\).  Substituting these bounds into \eqref{eq:Gamma_u_endpoint_exact} gives
\begin{equation}\label{eq:Gamma_u_lower_endpoint}
  \Gamma_u\ge
  2m_*\bigl(\alphastar(1-s^2)L(s)+\bstar(1-s^2)L(s(1+s^2))\bigr)-h(s^2)-L(s).
\end{equation}
Let
\begin{equation}\label{eq:delta_s}
  \delta(s)=L(s)-L(s(1+s^2)).
\end{equation}
Then
\begin{equation}\label{eq:Gamma_u_delta}
  \Gamma_u\ge
  \bigl(2m_*(1-s^2)-1\bigr)L(s)-2m_*\bstar(1-s^2)\delta(s)-h(s^2).
\end{equation}
For \(s\le1/25\),
\begin{equation}\label{eq:L_s_lower}
  L(s)\ge\log\frac1s-\frac1{24},
\end{equation}
while
\begin{equation}\label{eq:delta_bound}
  \delta(s)\le2s^2,
\end{equation}
and
\begin{equation}\label{eq:h_s2_bound}
  h(s^2)\le h(s_1^2).
\end{equation}
Set
\begin{equation}\label{eq:A0_C0}
  A_0=2m_*(1-s_1^2)-1,
  \qquad
  C_0=\frac{A_0}{24}+4m_*\bstar s_1^2+h(s_1^2).
\end{equation}
Numerically,
\begin{equation}\label{eq:A0C0num}
  A_0=0.232606493243870\ldots,
  \qquad
  C_0=0.0219863330048360\ldots .
\end{equation}
Thus
\begin{equation}\label{eq:Gamma_u_A0}
  \Gamma_u\ge A_0\log\frac1s-C_0.
\end{equation}
Finally,
\begin{equation}\label{eq:endpoint_margin}
  \left(A_0-\frac q2\right)\log25-C_0=0.349200203430100\ldots>0.
\end{equation}
Since \(\log(1/s)\ge\log25\), \eqref{eq:endpoint_positive} follows.
\end{proof}

\begin{proposition}[Certified branch structure]\label{prop:branch}
In the feasible region \eqref{eq:feasible_triangle}, the zero set \(F(u,v)=0\) consists of exactly two smooth arcs meeting at a single fold point.  More precisely:
\begin{enumerate}[label=(\roman*)]
\item The lower arc begins at the bifurcation point \((x_0,0)\), with
\begin{equation}\label{eq:x0_num}
  x_0=0.220715379651680\ldots,
\end{equation}
and ends at the fold
\begin{equation}\label{eq:fold_num}
  (u_f,v_f)=(0.552276058360800\ldots,0.434687834510229\ldots).
\end{equation}
Along this arc,
\begin{equation}\label{eq:lower_signs}
  \Gamma_{vv}<0,
  \qquad
  \Delta<0.
\end{equation}
\item The upper arc begins at the boundary point \((x,y)=(1,0)\) and ends at \((u_f,v_f)\).  On the compact part \(y\ge0.04\),
\begin{equation}\label{eq:upper_signs}
  \Gamma_{vv}>0,
  \qquad
  \Delta<0.
\end{equation}
For \(0<y\le0.04\), Lemma \ref{lem:endpoint} gives \(\Gamma_u>0\).
\item The endpoint values satisfy
\begin{equation}\label{eq:endpoint_values_branch}
  \Gamma_u(x_0,0)>0,
  \qquad
  \Gamma_u(u_f,v_f)>0.41.
\end{equation}
\end{enumerate}
\end{proposition}

\begin{proof}
The proposition is the certified interval statement proved in Appendix \ref{app:certificates}.  The proof uses one-variable interval Newton for the bifurcation point, interval Newton for the fold, Taylor-model interval arithmetic for \eqref{eq:lower_signs} and \eqref{eq:upper_signs}, and a uniform band certificate for the approach to \((1,0)\).
\end{proof}

\begin{lemma}[No off-diagonal critical points]\label{lem:no_offdiag}
There is no point \((x,y)\in(0,1)^2\) with \(x\ne y\) such that
\begin{equation}\label{eq:no_offdiag_condition}
  G_x(x,y)=G_y(x,y)=0.
\end{equation}
\end{lemma}

\begin{proof}
By symmetry, assume \(v>0\).  Any such critical point must satisfy \(F=0\) and \(\Gamma_u=0\).  On the lower arc, Proposition \ref{prop:branch} gives \(\Gamma_{vv}<0\) and \(\Delta<0\).  Hence by \eqref{eq:schur}, \(\Gamma_u\) is strictly increasing along the lower arc.  Since \(\Gamma_u(x_0,0)>0\), \(\Gamma_u>0\) throughout the lower arc.

On the upper arc with \(y\le0.04\), Lemma \ref{lem:endpoint} gives \(\Gamma_u>0\).  On the compact part of the upper arc, Proposition \ref{prop:branch} gives \(\Gamma_{vv}>0\) and \(\Delta<0\).  Hence by \eqref{eq:schur}, \(\Gamma_u\) decreases toward the fold.  Since \(\Gamma_u(u_f,v_f)>0.41\), it remains positive on this compact part.  Therefore no point on \(F=0\), \(v>0\), can satisfy \(\Gamma_u=0\).
\end{proof}

\section{Proof of the main inequality}\label{sec:mainineq}

\begin{theorem}[Bivariate entropy inequality]\label{thm:bivariate}
For every \(x,y\in[0,1]\),
\begin{equation}\label{eq:bivar_theorem}
  \mstar\bigl((1-\bstar)h(xy)+\bstar h(I(x,y))\bigr)\ge \frac{x h(y)+y h(x)}2.
\end{equation}
Equality occurs on the axes, at \((\xstar,\xstar)\), and at \((1,1)\).
\end{theorem}

\begin{proof}
This is exactly the assertion \(G(x,y)\ge0\).  Since \(G\) is continuous on \([0,1]^2\), it attains a minimum.  If a minimizer lies on the boundary, Section \ref{sec:boundary} gives nonnegativity.  If it lies in the interior and on the diagonal, Lemma \ref{lem:diagonal} gives nonnegativity.  If it lies in the interior and off the diagonal, it would be an off-diagonal critical point, contradicting Lemma \ref{lem:no_offdiag}.  Thus \(G\ge0\) everywhere.  The equality cases are the boundary equalities from Section \ref{sec:boundary} and the diagonal equality \((\xstar,\xstar)\) from Lemma \ref{lem:diagonal}.
\end{proof}

\section{The remaining distributional gap}\label{sec:gap}

We record explicitly why Theorem \ref{thm:bivariate} does not, by itself,
prove the union-closed consequence suggested by Liu's constant.

Let \(A\) be uniform on a union-closed family \(\Fcal\), and reveal coordinates
in order.  Let \(X_i\) denote the conditional probability that coordinate \(i\)
is absent, given the previous coordinates.  The chain rule gives
\begin{equation}\label{eq:gap_chain}
  H(A)=\sum_i \E h(X_i).
\end{equation}
For the independent two-sample coupling, if two histories give conditional
absence probabilities \(x\) and \(y\), the absence probability in the union is
\(xy\).  Thus the independent entropy contribution involves an expectation of
\(h(XY)\) with respect to a product distribution of conditional probabilities.

For a conditionally i.i.d. coupling, however, the entropy term involving
\begin{equation}
  I(x,y)=xy(1+(1-x)(1-y))
\end{equation}
is not generally averaged with respect to the same product distribution.  In
Liu's notation, the relevant optimization problem contains two numerator
terms: one averaged with respect to a product measure of the marginal mixture,
and one averaged with respect to a mixture of product measures.  Schematically,
if
\begin{equation}
  P=(1-q)P_0+qP_1,
\end{equation}
then the two terms have the form
\begin{equation}\label{eq:two_measures_gap}
  \E_{P^{\otimes2}} h(XY)
  \qquad\text{and}\qquad
  \E_{(1-q)P_0^{\otimes2}+qP_1^{\otimes2}} h(I(X,Y)).
\end{equation}
These are generally expectations with respect to different distributions on
\((X,Y)\).

The proof of Theorem \ref{thm:bivariate} establishes the pointwise inequality
\begin{equation}\label{eq:pointwise_gap}
  \mstar\bigl((1-\bstar)h(xy)+\bstar h(I(x,y))\bigr)
  \ge \frac{x h(y)+y h(x)}2.
\end{equation}
Averaging \eqref{eq:pointwise_gap} over a single joint distribution of
\((X,Y)\) is legitimate, but it does not imply a bound for the two-measure
quantity in \eqref{eq:two_measures_gap}.  In particular, the factorization
\begin{equation}\label{eq:wrong_factorization_gap}
  \E\frac{Xh(Y)+Yh(X)}2 = \E X\,\E h(X)
\end{equation}
requires an independence relation that is not available for the
conditionally i.i.d. term.

Consequently, the present paper should be read as a certified proof of the
bivariate inequality \eqref{eq:pointwise_gap}, together with a precise
identification of the remaining obstacle.  To obtain Liu's full
\(0.382709\ldots\) union-closed bound from this approach, one would need either
\begin{enumerate}
  \item a distributional inequality matching the two expectations in
  \eqref{eq:two_measures_gap}, or
  \item an independent proof that the worst case reduces to the same measure
  for both terms, such as a proof of the relevant optimizer-structure
  statement.
\end{enumerate}
We do not prove either statement here.

\appendix
\section{Computational certificates}\label{app:certificates}

This appendix records the finite certificates used in Sections
\ref{sec:diagonal} and \ref{sec:offdiagonal}.  All interval computations use
outward rounding and \(50\) decimal digits of precision.  The implementation is
available at
\begin{equation}\label{eq:code_url}
  \text{\url{https://github.com/kaare85-eng/liu-entropy-bound-certificate}}.
\end{equation}

\begin{proposition}[Constant enclosures]\label{prop:A1}
The constants \(\xstar\), \(\mstar\), \(\bstar\), \(x_0\), \(u_f\), and \(v_f\)
are enclosed in intervals of radius less than \(10^{-14}\) around the values
displayed in \eqref{eq:mstar_numeric}, \eqref{eq:betastar_numeric},
\eqref{eq:x0_num}, and \eqref{eq:fold_num}.  The enclosure for \(\xstar\)
certifies the unique root of \eqref{eq:xstarpoly} in \((0,1)\), and the other
constants are evaluated from their defining formulas.
\end{proposition}

\begin{proposition}[Bernstein positivity]\label{prop:A2}
For
\[
  P=P_0+\mstar P_1+\mstar\bstar P_2,
\]
all \(28\) degree-\(27\) Bernstein coefficients of \(P\) on \([0,1]\) are
positive.  The minimum certified lower bound is
\begin{equation}\label{eq:min_bern_app}
  \min_i b_i>4.550899799454681.
\end{equation}
Consequently,
\[
  P(x)>0\qquad(0\le x\le1).
\]
\end{proposition}

\begin{proposition}[Diagonal signs]\label{prop:A3}
The interval evaluation certifies
\begin{align}\label{eq:A3_signs}
&M(\mstar)>0,
\quad M'(\mstar)<0,
\quad M''(\mstar)>0,
\quad M^{(3)}(\mstar)>0,
\quad M^{(4)}(\mstar)<0,
\notag\\
&M''(\xstar)>0,
\quad M''(0.1)<0.
\end{align}
The numerical margins are those displayed in \eqref{eq:mstar_signs_num},
\eqref{eq:M2xstar}, and \eqref{eq:M2_01}.
\end{proposition}

\begin{proposition}[Unique bifurcation]\label{prop:A4}
Let
\begin{equation}\label{eq:Hdef_appendix}
  H(u)=\frac12\Gamma_{vv}(u,0).
\end{equation}
Then
\begin{equation}\label{eq:Hexplicit_appendix}
  H(u)
  =
  L(u)+\frac{1}{2(1-u)}
  -\mstar\alphastar L(u^2)
  -2\mstar\bstar(1-u+u^2)L(D(u)).
\end{equation}
The interval proof certifies
\[
  H(u)<0\qquad(0<u<x_0),
\]
\[
  H(u)>0\qquad(x_0<u<1),
\]
and a unique zero at \(x_0\).  On the root box, \(H'(u)>0\).
\end{proposition}

\begin{proposition}[Fold uniqueness and value]\label{prop:A5}
The system
\begin{equation}\label{eq:fold_system_app}
  F(u,v)=0,
  \qquad
  \Gamma_{vv}(u,v)=0
\end{equation}
has a unique solution in the fold box centered at \((u_f,v_f)\).  At this
solution,
\begin{equation}\label{eq:fold_Gammau_app}
  \Gamma_u(u_f,v_f)>0.41.
\end{equation}
Moreover, the interval evaluation gives
\begin{equation}\label{eq:fold_delta_app}
  \Delta(u_f,v_f)<-31.17.
\end{equation}
\end{proposition}

\paragraph{Proof of branch isolation.}
We now explain how the preceding certificates are assembled into the branch
isolation statement.  Consider the zero set
\[
  \{(u,v):F(u,v)=0,\ v>0\}
\]
inside the feasible triangle
\[
  0<u<1,
  \qquad
  0<v<\min(u,1-u).
\]
Any connected component of this zero set must fall into one of the following
topological alternatives.

First, it may be a closed loop in the interior.  Such a loop must have at
least two extrema in the \(u\)-projection, and hence at least two fold points,
that is, at least two solutions of
\[
  F(u,v)=0,
  \qquad
  F_v(u,v)=0.
\]
On \(F=0\) this is equivalent to
\[
  F(u,v)=0,
  \qquad
  \Gamma_{vv}(u,v)=0.
\]
Proposition \ref{prop:A5} certifies that there is only one such fold.  Thus
closed loops are excluded.

Second, a component may meet the diagonal \(v=0\).  Since \(\Gamma\) is even in
\(v\), branches can bifurcate from the diagonal only where
\[
  \Gamma_{vv}(u,0)=0.
\]
Proposition \ref{prop:A4} certifies that the only such point is \(u=x_0\).
Thus the only possible diagonal birth is the already traced lower branch.

Third, a component may approach the boundary of the feasible triangle away
from the corner \((x,y)=(1,0)\).  The side \(y=0\), away from \(x=1\), is
excluded analytically: for fixed \(0<x<1\),
\[
  F(x,y)\to -\infty
  \qquad (y\to0^+).
\]
The side \(x=1\), away from \(y=0\), is also excluded analytically: for fixed
\(0<y<1\),
\[
  F(x,y)\to +\infty
  \qquad (x\to1^-).
\]
The reflected sides \(x=0\) and \(y=1\) do not occur as distinct sides in the
half-triangle \(v>0\); they are obtained by symmetry.

Fourth, a component may approach the corner \((x,y)=(1,0)\).  The upper arc
has exactly this endpoint.  The uniqueness of this approach is certified by a
one-dimensional interval Newton band computation in the variable
\(s=y\).  For \(s\in[0.04,0.08]\), the computation certifies a unique solution
of \(F(\varepsilon,s)=0\), where \(\varepsilon=1-x\), in the upper-branch tube.
The eight \(s\)-bands have width \(0.005\); the maximal Newton contraction
factor is \(0.6296\).  For \(0<s<0.04\), Lemma \ref{lem:endpoint} gives the
analytic endpoint estimate and identifies the same unique branch approaching
the corner.

The two known arcs account for the unique diagonal bifurcation, the unique
fold, and the unique corner approach.  The alternatives above exhaust all
possible connected components.  Therefore there are no additional components
of \(F=0\), \(v>0\), in the feasible triangle.

\begin{proposition}[Branch isolation]\label{prop:A6}
The zero set
\[
  F(u,v)=0,
  \qquad
  v>0,
\]
in the feasible triangle \eqref{eq:feasible_triangle} consists of exactly the
two arcs stated in Proposition \ref{prop:branch}.  The lower arc runs from the
diagonal bifurcation point \((x_0,0)\) to the fold \((u_f,v_f)\).  The upper
arc runs from the endpoint \((x,y)=(1,0)\) to the same fold.
\end{proposition}

\begin{proposition}[Lower arm determinant sign]\label{prop:A7}
On the lower arc of \(F=0\), the Taylor-model interval certificate verifies
\begin{equation}\label{eq:lower_det_app}
  \Gamma_{vv}<0,
  \qquad
  \Delta<0.
\end{equation}
The bifurcation zone uses \((u,w)\)-coordinates with \(w=v^2\).  The
Taylor-model cover contains \(23\) boxes in the near-bifurcation zone and
\(56\) boxes in the compact lower zone.  The largest certified upper bound for
\(\Delta\) is
\begin{equation}\label{eq:lower_margin_app}
  -0.00116054425855.
\end{equation}
\end{proposition}

\begin{proposition}[Upper arm determinant sign]\label{prop:A8}
On the compact upper arc with \(y\in[0.08,y_f]\), Taylor-model interval
arithmetic in the coordinates
\[
  (\eps,s)=(1-x,y)
\]
verifies
\begin{equation}\label{eq:upper_det_app}
  \Gamma_{vv}>0,
  \qquad
  \Delta<0.
\end{equation}
The cover uses \(42\) adaptive boxes with maximum step size at most \(0.001\).
The largest certified upper bound for \(\Delta\) is
\begin{equation}\label{eq:upper_margin_app}
  -31.1897707236,
\end{equation}
and the smallest certified lower bound for \(\Gamma_{vv}\) is
\begin{equation}\label{eq:upper_gvv_margin_app}
  0.000394450122601.
\end{equation}
\end{proposition}

\begin{proposition}[Endpoint positivity]\label{prop:A9}
For the endpoint region \(0<y\le0.04\), Lemma \ref{lem:endpoint} gives the
analytic lower bound
\begin{equation}\label{eq:A9_endpoint}
  \Gamma_u\ge\frac{2\mstar-1}{2}\log\frac1y>0.
\end{equation}
At the transition point \(y=0.04\), the computed value is
\begin{equation}\label{eq:endpoint_numeric_app}
  \Gamma_u=0.745513582988373\ldots,
  \qquad
  \frac{2\mstar-1}{2}\log25=0.377544881375125\ldots .
\end{equation}
\end{proposition}


\section*{Acknowledgements}
This paper was developed with substantial AI assistance.
The mathematical strategy, proof structure, and all
certificate computations were carried out through
extended interaction with large language models.

\begin{thebibliography}{99}

\bibitem{Frankl1979}
P.~Frankl.
\newblock Union-closed sets conjecture.
\newblock Problem statement, 1979. See also D. B. West's \emph{Open Problems in Graph Theory and Combinatorics}.

\bibitem{BruhnSchaudt2015}
H.~Bruhn and O.~Schaudt.
\newblock The journey of the union-closed sets conjecture.
\newblock \emph{Graphs and Combinatorics}, 31(6):2043--2074, 2015.

\bibitem{Gilmer2022}
J.~Gilmer.
\newblock A constant lower bound for the union-closed sets conjecture.
\newblock arXiv:2211.09055, 2022.

\bibitem{AlweissHuangSellke2024}
R.~Alweiss, B.~Huang, and M.~Sellke.
\newblock Improved lower bound for Frankl's union-closed sets conjecture.
\newblock \emph{The Electronic Journal of Combinatorics}, 31(3):P3.35, 2024.

\bibitem{ChaseLovett2022}
Z.~Chase and S.~Lovett.
\newblock Approximate union closed conjecture.
\newblock arXiv:2211.11689, 2022.

\bibitem{Sawin2022}
W.~Sawin.
\newblock An improved lower bound for the union-closed sets conjecture.
\newblock arXiv:2211.11504, 2022.

\bibitem{Liu2023}
J.~Liu.
\newblock Improving the lower bound for the union-closed sets conjecture via conditionally {IID} coupling.
\newblock arXiv:2306.08824, 2023.

\bibitem{Cambie2022}
S.~Cambie.
\newblock Better bounds for the union-closed sets conjecture using the entropy approach.
\newblock arXiv:2212.12500, 2022.

\bibitem{Yu2023}
L.~Yu.
\newblock Dimension-free bounds for the union-closed sets conjecture.
\newblock \emph{Entropy}, 25(5):767, 2023.


\bibitem{Pebody2022}
L.~Pebody.
\newblock Extension of a method of Gilmer.
\newblock arXiv:2211.13139, 2022.

\bibitem{Boppana1985}
R.~B. Boppana.
\newblock Amplification of probabilistic Boolean formulas.
\newblock In \emph{26th Annual Symposium on Foundations of Computer Science (SFCS 1985)}, pages 20--29. IEEE, 1985.

\bibitem{PulajWood2024}
J.~Pulaj and K.~Wood.
\newblock Local configurations in union-closed families.
\newblock \emph{Experimental Mathematics}, to appear; arXiv:2301.01331, 2024.

\end{thebibliography}


\end{document}
